% Options for packages loaded elsewhere
\PassOptionsToPackage{unicode}{hyperref}
\PassOptionsToPackage{hyphens}{url}
%
\documentclass[
  doc,floatsintext]{apa6}
\usepackage{amsmath,amssymb}
\usepackage{iftex}
\ifPDFTeX
  \usepackage[T1]{fontenc}
  \usepackage[utf8]{inputenc}
  \usepackage{textcomp} % provide euro and other symbols
\else % if luatex or xetex
  \usepackage{unicode-math} % this also loads fontspec
  \defaultfontfeatures{Scale=MatchLowercase}
  \defaultfontfeatures[\rmfamily]{Ligatures=TeX,Scale=1}
\fi
\usepackage{lmodern}
\ifPDFTeX\else
  % xetex/luatex font selection
\fi
% Use upquote if available, for straight quotes in verbatim environments
\IfFileExists{upquote.sty}{\usepackage{upquote}}{}
\IfFileExists{microtype.sty}{% use microtype if available
  \usepackage[]{microtype}
  \UseMicrotypeSet[protrusion]{basicmath} % disable protrusion for tt fonts
}{}
\makeatletter
\@ifundefined{KOMAClassName}{% if non-KOMA class
  \IfFileExists{parskip.sty}{%
    \usepackage{parskip}
  }{% else
    \setlength{\parindent}{0pt}
    \setlength{\parskip}{6pt plus 2pt minus 1pt}}
}{% if KOMA class
  \KOMAoptions{parskip=half}}
\makeatother
\usepackage{xcolor}
\usepackage{graphicx}
\makeatletter
\def\maxwidth{\ifdim\Gin@nat@width>\linewidth\linewidth\else\Gin@nat@width\fi}
\def\maxheight{\ifdim\Gin@nat@height>\textheight\textheight\else\Gin@nat@height\fi}
\makeatother
% Scale images if necessary, so that they will not overflow the page
% margins by default, and it is still possible to overwrite the defaults
% using explicit options in \includegraphics[width, height, ...]{}
\setkeys{Gin}{width=\maxwidth,height=\maxheight,keepaspectratio}
% Set default figure placement to htbp
\makeatletter
\def\fps@figure{htbp}
\makeatother
\setlength{\emergencystretch}{3em} % prevent overfull lines
\providecommand{\tightlist}{%
  \setlength{\itemsep}{0pt}\setlength{\parskip}{0pt}}
\setcounter{secnumdepth}{-\maxdimen} % remove section numbering
% Make \paragraph and \subparagraph free-standing
\makeatletter
\ifx\paragraph\undefined\else
  \let\oldparagraph\paragraph
  \renewcommand{\paragraph}{
    \@ifstar
      \xxxParagraphStar
      \xxxParagraphNoStar
  }
  \newcommand{\xxxParagraphStar}[1]{\oldparagraph*{#1}\mbox{}}
  \newcommand{\xxxParagraphNoStar}[1]{\oldparagraph{#1}\mbox{}}
\fi
\ifx\subparagraph\undefined\else
  \let\oldsubparagraph\subparagraph
  \renewcommand{\subparagraph}{
    \@ifstar
      \xxxSubParagraphStar
      \xxxSubParagraphNoStar
  }
  \newcommand{\xxxSubParagraphStar}[1]{\oldsubparagraph*{#1}\mbox{}}
  \newcommand{\xxxSubParagraphNoStar}[1]{\oldsubparagraph{#1}\mbox{}}
\fi
\makeatother
% definitions for citeproc citations
\NewDocumentCommand\citeproctext{}{}
\NewDocumentCommand\citeproc{mm}{%
  \begingroup\def\citeproctext{#2}\cite{#1}\endgroup}
\makeatletter
 % allow citations to break across lines
 \let\@cite@ofmt\@firstofone
 % avoid brackets around text for \cite:
 \def\@biblabel#1{}
 \def\@cite#1#2{{#1\if@tempswa , #2\fi}}
\makeatother
\newlength{\cslhangindent}
\setlength{\cslhangindent}{1.5em}
\newlength{\csllabelwidth}
\setlength{\csllabelwidth}{3em}
\newenvironment{CSLReferences}[2] % #1 hanging-indent, #2 entry-spacing
 {\begin{list}{}{%
  \setlength{\itemindent}{0pt}
  \setlength{\leftmargin}{0pt}
  \setlength{\parsep}{0pt}
  % turn on hanging indent if param 1 is 1
  \ifodd #1
   \setlength{\leftmargin}{\cslhangindent}
   \setlength{\itemindent}{-1\cslhangindent}
  \fi
  % set entry spacing
  \setlength{\itemsep}{#2\baselineskip}}}
 {\end{list}}
\usepackage{calc}
\newcommand{\CSLBlock}[1]{\hfill\break\parbox[t]{\linewidth}{\strut\ignorespaces#1\strut}}
\newcommand{\CSLLeftMargin}[1]{\parbox[t]{\csllabelwidth}{\strut#1\strut}}
\newcommand{\CSLRightInline}[1]{\parbox[t]{\linewidth - \csllabelwidth}{\strut#1\strut}}
\newcommand{\CSLIndent}[1]{\hspace{\cslhangindent}#1}
\ifLuaTeX
\usepackage[bidi=basic]{babel}
\else
\usepackage[bidi=default]{babel}
\fi
\babelprovide[main,import]{english}
% get rid of language-specific shorthands (see #6817):
\let\LanguageShortHands\languageshorthands
\def\languageshorthands#1{}
% Manuscript styling
\usepackage{upgreek}
\captionsetup{font=singlespacing,justification=justified}

% Table formatting
\usepackage{longtable}
\usepackage{lscape}
% \usepackage[counterclockwise]{rotating}   % Landscape page setup for large tables
\usepackage{multirow}		% Table styling
\usepackage{tabularx}		% Control Column width
\usepackage[flushleft]{threeparttable}	% Allows for three part tables with a specified notes section
\usepackage{threeparttablex}            % Lets threeparttable work with longtable

% Create new environments so endfloat can handle them
% \newenvironment{ltable}
%   {\begin{landscape}\centering\begin{threeparttable}}
%   {\end{threeparttable}\end{landscape}}
\newenvironment{lltable}{\begin{landscape}\centering\begin{ThreePartTable}}{\end{ThreePartTable}\end{landscape}}

% Enables adjusting longtable caption width to table width
% Solution found at http://golatex.de/longtable-mit-caption-so-breit-wie-die-tabelle-t15767.html
\makeatletter
\newcommand\LastLTentrywidth{1em}
\newlength\longtablewidth
\setlength{\longtablewidth}{1in}
\newcommand{\getlongtablewidth}{\begingroup \ifcsname LT@\roman{LT@tables}\endcsname \global\longtablewidth=0pt \renewcommand{\LT@entry}[2]{\global\advance\longtablewidth by ##2\relax\gdef\LastLTentrywidth{##2}}\@nameuse{LT@\roman{LT@tables}} \fi \endgroup}

% \setlength{\parindent}{0.5in}
% \setlength{\parskip}{0pt plus 0pt minus 0pt}

% Overwrite redefinition of paragraph and subparagraph by the default LaTeX template
% See https://github.com/crsh/papaja/issues/292
\makeatletter
\renewcommand{\paragraph}{\@startsection{paragraph}{4}{\parindent}%
  {0\baselineskip \@plus 0.2ex \@minus 0.2ex}%
  {-1em}%
  {\normalfont\normalsize\bfseries\itshape\typesectitle}}

\renewcommand{\subparagraph}[1]{\@startsection{subparagraph}{5}{1em}%
  {0\baselineskip \@plus 0.2ex \@minus 0.2ex}%
  {-\z@\relax}%
  {\normalfont\normalsize\itshape\hspace{\parindent}{#1}\textit{\addperi}}{\relax}}
\makeatother

\makeatletter
\usepackage{etoolbox}
\patchcmd{\maketitle}
  {\section{\normalfont\normalsize\abstractname}}
  {\section*{\normalfont\normalsize\abstractname}}
  {}{\typeout{Failed to patch abstract.}}
\patchcmd{\maketitle}
  {\section{\protect\normalfont{\@title}}}
  {\section*{\protect\normalfont{\@title}}}
  {}{\typeout{Failed to patch title.}}
\makeatother

\usepackage{xpatch}
\makeatletter
\xapptocmd\appendix
  {\xapptocmd\section
    {\addcontentsline{toc}{section}{\appendixname\ifoneappendix\else~\theappendix\fi: #1}}
    {}{\InnerPatchFailed}%
  }
{}{\PatchFailed}
\makeatother
\keywords{Trust in science; science knowledge; deficit model}
\usepackage{csquotes}
\usepackage{placeins} 

\ifLuaTeX
  \usepackage{selnolig}  % disable illegal ligatures
\fi
\usepackage{bookmark}
\IfFileExists{xurl.sty}{\usepackage{xurl}}{} % add URL line breaks if available
\urlstyle{same}
\hypersetup{
  pdftitle={The rational impression account of trust in science},
  pdflang={en-EN},
  pdfkeywords={Trust in science; science knowledge; deficit model},
  hidelinks,
  pdfcreator={LaTeX via pandoc}}

\title{The rational impression account of trust in science}
\author{\textsuperscript{}}
\date{}


\shorttitle{Rational impression account}

\affiliation{\vspace{0.5cm}\textsuperscript{} }

\abstract{%
Trust in science matters to a range of significant outcomes such as intent to vaccination and believe in climate change. Around the globe, most people trust science at least to some extent. However, the causes of this trust aren't well understood. Here we test a rational impression account of trust in science: According to this account, people trust scientists because they are impressed by their findings. This impression persists even after knowledge on the specific content of the findings has vanished. We present evidence in line with the preditions in two Experiments (total n = XX). In Experiment 1, we show that impressive scientific findings lead people to think of the scientists as more competent and their discipline as more trustworthy. In Experiment 2, we find that these impressions are partly independent of remembering specific findings. The rational impression account could explain the following stylized observations: (i) Globally, most people have at least some trust in science, and almost everyone, at least in the US, trusts (nearly all of) basic science; (ii) Education, and in particular science education, is the best predictor of trust in science; (iii) The levels of specific science knowledge are relatively low, as is the associations between science knowledge and trust in science.
}



\begin{document}
\maketitle

\section{Introduction}\label{introduction}

Jon D. Miller, a pioneer in the field of public understanding of science, found that at the end of the 20th century, approximately 17\% of Americans qualified as scientifically literate--a measure he described as the ``level of skill required to read most of the articles in the Tuesday science section of The New York Times''. In the early 1990s, his estimate had been even lower, around 10\%. Meanwhile, the level of trust in the US has been relatively high and remarkably stable (Funk \& Kennedy, 2020; Funk, Tyson, Kennedy, \& Johnson, 2020; Smith \& Son, 2013): Since the 1970s, on average 43.1\% of Americans say they have a great deal of confidence in the scientific community, the second highest score among 13 institutions inclueded in the US General Social Survey (GSS), just behind the military (see Fig. \ref{fig:gss-trend}). How can we explain this gap between understanding science and trusting it? Here, we argue that people might mostly trust science because they have been impressed by it.



\begin{figure}
\centering
\includegraphics{output/figures/gss-trend.pdf}
\caption{\label{fig:gss-trend}\emph{Top four trusted institutions in the US}. The plot shows the variation of confidence in the top four trusted institutions in the US over time. The data for this plot is from the cumulative file of the US General Social Survey (GSS).}
\end{figure}

We begin by presenting global evidence that most people trust science at least to some extent. We then review existing explanations for why people trust science, in particular the deficit-model, and argue that several empirical observations violate their predictions. We proceed by proposing an alternative model, which we term the `rational impression account'. In two Experiments, we provide evidence in favour of this account.

Trust in science matters. It is related to many desirable outcomes, from acceptance of anthropogenic climate change (Cologna \& Siegrist, 2020) or vaccination (Lindholt, Jørgensen, Bor, \& Petersen, 2021; Sturgis, Brunton-Smith, \& Jackson, 2021) to following recommendations during COVID (Algan, Cohen, Davoine, Foucault, \& Stantcheva, 2021).

Across the globe, most people report to trust science at least to some extent. In 2018, the Wellcome Global Monitor collected responses from more than 140,000 people in over 140 countries on their trust in science (Wellcome Global Monitor, 2018). Globally, 34\% said they trust science `a lot', 45\% answered `some' and only 13\% chose the option `not much, or not at all' (10\% said they did not know). In 2020, a follow-up survey on 119.000 people from 113 countries found that in midst of the Covid-19 pandemic and before vaccines were available, trust in science and scientists was even higher: 41\% of respondents said they trust science `a lot', 39\% `some', and 13\% `not much, or not at all' (Wellcome Global Monitor, 2020). In late 2022 and early 2023, another large survey project collected data on trust in science from 71,922 respondents in 68 countries. They found trust in scientists to be ``moderately high'' (mean = 3.62; sd= 0.70; Scale: 1 = very low, 2 = somewhat low, 3 = neither high nor low, 4 = somewhat high, 5 = very high). No country showed low overall trust in scientists. A recent study in the US found that even participants who say they do not trust science and who believe anti-science conspiracy theories (e.g.~that the earth is flat) did in fact accept the scientific consensus on basic, non-politicized knowledge questions (e.g.~`Are electrons are smaller than atoms?') in most cases (Pfänder, Kerzreho, \& Mercier, 2024).

How can we explain trust in science? The literature on the public understanding of science has generally focused not on explaining trust, but a lack of it, and more generally negative attitudes towards science (Bauer, Allum, \& Miller, 2007).

One explanation, is what Gauchat (2011) coined the \emph{alienation model}. According to the alienation model, the ``public disassociation with science is a symptom of a general disenchantment with late modernity, mainly, the limitations associated with codified expertise, rational bureaucracy, and institutional authority'' (Gauchat, 2011). This explanation builds on the work of social theorists (Beck, 1992; Giddens, 1991; Habermas, 1989; see Gauchat, 2011 for an overview). According to these theorists, a modern, complex world increasingly requires expertise, and thus shapes institutions of knowledge elites. People who are not part of these institutions, experience a lack of agency, resulting in a feeling of alienation.

The alienation model would predict increasing distrust in science. This does not sit well with the observation depicted in Fig. \ref{fig:gss-trend}: At least in the US, where long-term data on this is available, trust in science appears to be mostly stable, on average (Funk et al., 2020; see also Gauchat, 2012).

A second explanation which has dominated the literature on public understanding of science for decades (Bauer et al., 2007) is widely known as he ``deficit model''. According to the deficit model, people do not trust science enough, because they do not know enough about it. The deficit model implies that when people trust science, it is because they know a lot about it (Bauer et al., 2007).

Seemingly, empirical observations are in line with that view: large scale survey studies consistently identify education, and in particular science education--among all sorts of demographic variables--to be the most strongly correlated with trust in science (Bak, 2001; Noy \& O'Brien, 2019; Wellcome Global Monitor, 2018, 2020): The more educated tend to trust science more. The deficit model provides a straightforward causal explanation for this correlation: education fosters trust by improving the understanding of science.

However, the evidence for this causal model is far from conclusive. To begin with, there is no consensus among researchers on how to measure understanding of science (Gauchat \& Andrews, 2018). An established, widely used measure consists of asking people a set of basic science knowledge questions (see e.g. Durant, Evans, \& Thomas, 1989; Miller, 1998). Sometimes it is also measured by coding answers to open-ended questions about scientific methods. By either measure past research has found science knowledge to be, in some cases, alarming low (Miller, 2004; National Academies of Sciences, Engineering, and Medicine, 2016), contrasting with the high levels of trust in science, and to be at best weakly associated with attitudes towards science (Allum, Sturgis, Tabourazi, \& Brunton-Smith, 2008). Other research has shown that the correlation between education and positive attitudes towards science holds even after controlling for science knowledge, which led the researchers to conclude that only part of the effect of education could be explained by transmitting science knowledge (Bak, 2001).

\section{The rational impression account}\label{the-rational-impression-account}

Here, we seek to explain the following stylized facts: (i) Globally, most people have at least some trust in science, and almost everyone, at least in the US, trusts (nearly all of) basic science; (ii) Education, and in particular science education, is the best predictor of trust in science; (iii) The levels of specific science knowledge are relatively low, as is the associations between science knowledge and trust in science.

All of these observations are compatible with a model of trust in science which we call the the `rational impression account'. According to this account, people are impressed by scientific findings, which in turn makes them trust scientists. The impression persists even after knowledge on the specific content of the findings has vanished.

For an information to be impressive, at least two criteria should be met: (i) it is perceived as hard to uncover and (ii) there is reason to believe is true. By (i), we mean information that most people would have no idea how to find out themselves. For example, most people would probably not be much impressed by someone telling them the biggest local park tree has exactly 110'201 leaves. Even though obtaining this information implies an exhausting counting effort, everyone knows how to do it. By contrast, finding out that it takes light \href{https://imagine.gsfc.nasa.gov/features/cosmic/milkyway_info.html}{approximately 100'000 years to travel from one end of the Milky Way to the other} is probably impressive to most people. Regarding (ii), if people cannot verify the information for themselves--as can be assumed to be the case for most scientific results--they need to rely on other cues to establish whether it might be true or not. One such cue is the degree of consensus between different sources. It has been shown that when they lack relevant background knowledge and have no reason to suspect a conspiracy, people take social consensus as a cue for accuracy of the information and competence of the informant (Pfänder, De Courson, \& Mercier, 2025). Science is perhaps the real-world case where this inference should matter most: If independent scientists agree on something like the length of the Milky Way, that should be impressive and make us think of the scientists as competent, even if we have no idea how they did it.

Once people are impressed, that impression tends to last, even if specific content is forgotten. We commonly form impressions of the people around us while forgetting the details of how we formed these impressions: If a colleague fixes our computer, we might forget exactly how they fixed it, yet remember that they are good at fixing computers. For an extreme example, it has been demonstrated that patients with severe amnesia continued to experience emotions linked to events they could not recall (Feinstein, Duff, \& Tranel, 2010). More related to our context, it has been shown that while people find some science-related explanations more satisfying than others, this hardly predicted how well they could recall the explanations shortly after (Liquin \& Lombrozo, 2022), suggesting that that impressions and knowledge can be quite detached.

The rational impression account of trust in science can reconcile the stylized facts outlined above: For (i), most people in the world receive at least some basic education that exposes them to impressive scientific findings, installing a baseline of trust. For (ii), education can be assumed to be the primary place of exposure to science, making it the most important variable. For (iii), specific knowledge about or understanding of science is not a requirement for trust.

While the rational impression account is plausible, there is no direct mechanistic test of it. In two Experiments, we aim to provide such as test.

\section{The present studies}\label{the-present-studies}

In this paper, we present two pre-registered Experiments (total n = 696) to test the predictions of the rational impression account. In Experiment 1, we test whether impressive scientific findings lead people to think of the scientists as more competent and trustworthy. In Experiment 2, we test whether that impression persists even if specific content is forgotten.

All experiments were preregistered, and choices of sample size were informed by power simulations. All materials, data and code can be found on Open Science Framework \href{https://osf.io/j3bk4/}{project page} (\url{https://osf.io/j3bk4/}). All analyses were conducted in R (version 4.2.2) using R Studio. We used two-sided tests even for directional hypotheses. Unless mentioned otherwise, we report unstandardized model coefficients that can be interpreted in units of the scales we use for our dependent variables.

For full transparency, we note that as part of this project, we conducted two additional Experiments which we present in detail in the ESM. The first one is almost identical to Experiment 1, but suffered from a minor technical error during the implementation. It replicates the findings of Experiment 1. The second one was supposed to test the same hypothesis as Experiment 2, but the treatment--a very short distraction task--did not alter any of the outcome variables, suggesting that our outcome measures were not sensitive enough.

\subsection{Experiment 1}\label{experiment-1}

The main goal of experiment one to see if exposure to impressive science content could enhance people's trust in scientists.

After providing their consent to participate in the study and passing an attention check, participants read the following introduction: ``You're going to read two texts about scientific discoveries, in two fields: archeology (the study of human activity through the recovery and analysis of material culture), and entomology (the study of insects). We ask you to read the texts, and then answer a few questions about them.'' Next, participants read vignettes about scientific findings in the disciplines of entomology and archaeology. We manipulated the impressiveness of the texts, so that there was one ``basic'' and one ``impressive'' version for each of the disciplines (see Table \ref{tab:exp1-stimuli}). We varied impressiveness within participants, but between disciplines: each participant saw an impressive version for one discipline, and a basic version for the other discipline. After reading each text, participants were asked a series of questions. We asked participants how much they thought they learned from reading the text (``How much do you feel you've learnt about human history by reading this text?'' {[}1 - Nothing, 2 - A bit, 3 - Some, 4 - Quite a bit, 5 - A lot{]}) and how impressive they found it (``How impressive do you think the findings of the archaeologists described in the text are?'' {[}1 - Not very impressive, 2 - A bit impressive, 3 - Quite impressive, 4 - Very impressive, 5 - Extremely impressive{]}). We also asked them whether reading the text changed their impression of the competence of the scientists of the respective discipline (``Would you agree that reading this text has made you think of archaeologists as more competent than you thought before?'' {[}1 - Strongly disagree, 2 - Disagree, 3 - Neither agree nor disagree, 4 - Agree, 5 - Strongly agree{]}) and their trust in the respective discipline (``Having read this text, would you agree that you trust the discipline of archaeology more than you did before?'' {[}1 - Strongly disagree, 2 - Disagree, 3 - Neither agree nor disagree, 4 - Agree, 5 - Strongly agree{]}). Finally, we also asked about the perceived consensus (``To which extent do you think the findings from the short text you just read reflect a minority or a majority opinion among archaeologists?'' {[}1 - Small minority, 2 - Minority, 3 - About half, 4 - Majority, 5 - Large majority{]}).

We had the following hypotheses:

\textbf{H1a: Participants will perceive scientists as more competent than they did before after having read an impressive text about their discipline's findings, compared to when reading a basic text.}

\textbf{H1b: Across all conditions, participants who are more impressed by the text about a discipline will also tend to perceive the scientists of the discipline as more competent.}

\textbf{H2a: Participants will trust a discipline more than they did before, after reading an impressive text about the discipline's findings, compared to when reading a basic text.}

\textbf{H2b: Across all conditions, participants who are more impressed by the text about a discipline will also tend to trust the scientists of the discipline more.}

We had the following research questions:

\textbf{RQ1: Do participants perceive that they learn more from the texts in the impressive condition, compared to the basic condition?}

\textbf{RQ2: Do perceptions of consensus interact with the relationships proposed in the hypotheses, such that greater perceived consensus is associated with a more positive relationship between impressiveness and trust/competence?}

\subsection{Methods}\label{methods}

\subsubsection{Participants}\label{participants}

A power simulation (see OSF) suggested that the minimum required sample size to detect a statistically significant effect for all hypotheses with a power of 0.9 is n = 100 participants. We therefore recruited 100 participants from UK via Prolific. One participant failed our attention check, resulting in a final sample of 99 participants (49 female, 50 male; \(age_\text{mean}\): 42.07, \(age_\text{sd}\): 12.86, \(age_\text{median}\): 40).

\subsubsection{Materials}\label{materials}

Table \ref{tab:exp1-stimuli} shows the vignettes we used.

\begingroup\fontsize{8}{10}\selectfont

\begin{longtable}[t]{>{\raggedright\arraybackslash}p{5em}>{\raggedright\arraybackslash}p{25em}>{\raggedright\arraybackslash}p{25em}}
\caption{\label{tab:exp1-stimuli}Stimuli of Experiment 1}\\
\toprule
 & Impressive & Basic\\
\midrule
\endfirsthead
\caption[]{\label{tab:exp1-stimuli}Stimuli of Experiment 1 \textit{(continued)}}\\
\toprule
 & Impressive & Basic\\
\midrule
\endhead

\endfoot
\bottomrule
\endlastfoot
Archeology & Archaeologists, scientists who study human history and prehistory, are able to tell, from their bones, whether someone was male or female, how old they were, and whether they suffered from a range of diseases. Archaeologists can now tell at what age someone, dead for tens of thousands of years, stopped drinking their mother’s milk, from the composition of their teeth.
Archaeologists learn about the language that our ancestors or cousins might have had. For instance, the nerve that is used to control breathing is larger in humans than in apes, plausibly because we need more fine-grained control of our breathing in order to speak. As a result, the canal containing that nerve is larger in humans than in apes – and it is also enlarged in Neanderthals.
Archaeologists can also tell, from an analysis of the tools they made, that most Neanderthals were right-handed. It’s thought that handedness is related to the evolution of language, another piece of evidence suggesting that Neanderthals likely possessed a form of language. & Archaeology is the science that studies human history and prehistory based on the analysis of objects from the past such as human bones, engravings, constructions, and various objects, from nails to bits of pots. This task requires a great deal of carefulness, because objects from the past need to often be dug out from the ground and patiently cleaned, without destroying them in the process.

Archaeologists have been able to shed light on human history in all continents, from ancient Egypt to the Incas in Peru or the Khmers in Cambodia.

Archaeologists study the paintings made by our ancestors, such as those that can be found in Lascaux, a set of caves found in the south of France that have been decorated by people at least 30000 years ago.

Archaeologists have also found remains of our more distant ancestors, showing that our species is just one among several that appeared, and then either changed or went extinct, such as Neanderthals, Homo erectus, or Homo habilis.\\
Entomology & Entomologists are the scientists who study insects. Some of them have specialized in understanding how insects perceive the world around them, and they have uncovered remarkable abilities. 

Entomologists interested in how flies’ visual perception works have used special displays to present images for much less than the blink of an eye, electrodes to record how individual cells in the flies’ brain react, and ultra-precise electron microscopy to examine their eyes. Thanks to these techniques, they have shown that some flies can perceive images that are displayed for just three milliseconds (a thousandth of a second) – about ten times shorter than a single movie frame (of which there are 24 per second). 

Entomologists who study the hair of crickets have shown that these microscopic hairs, which can be found on antenna-like organs attached to the crickets’ rear, are maybe the most sensitive organs in the animal kingdom. The researchers used extremely precise techniques to measure how the hair reacts to stimuli, such as laser-Doppler velocimetry, a technique capable of detecting the most minute of movements. They were able to show that the hair could react to changes in the motion of the air that had less energy than one particle of light, a single photon. & Entomologists are scientists who investigate insects, typically having a background in biology. They study, for example, how a swarm of bees organizes, or how ants communicate with each other. 

They also study how different insects interact with each other and their environment, whether some species are in danger of going extinct, or whether others are invasive species that need to be controlled.

Sometimes entomologists study insects by observing them in the wild, sometimes they conduct controlled experiments in laboratories, to see for example how different environmental factors change the behavior of insects, or to track exactly the same insects over a longer period of time.

An entomologist often specializes in one type of insect in order to study it in depth. For example, an entomologist who specializes in ants is called a myrmecologist.\\*
\end{longtable}
\endgroup{}

\subsection{Results and discussion}\label{results-and-discussion}

For the manipulation check, we find that the impressive text as more impressive (mean = 4.01, sd = 0.87; \(\hat{b}\) = 0.81 {[}0.569, 1.048{]}, p \textless{} .001) than the basic text (mean = 3.20, sd = 1.13). This overall difference was reflected in both the vignettes on entomology (\(\hat{b}\) = 0.94 {[}0.543, 1.339{]}, p \textless{} .001; \(\text{mean}_{\text{basic}}\) = 3.10, \(\text{sd}_{\text{basic}}\) = 1.09, \(\text{mean}_{\text{impressive}}\) = 4.04, \(\text{sd}_{\text{impressive}}\) = 0.89) and the vignettes on archaeology (\(\hat{b}\) = 0.67 {[}0.262, 1.086{]}, p = 0.002; \(\text{mean}_{\text{basic}}\) = 3.31, \(\text{sd}_{\text{basic}}\) = 1.18, \(\text{mean}_{\text{impressive}}\) = 3.98, \(\text{sd}_{\text{impressive}}\) = 0.87).

Participants perceived scientists as more competent after having read an impressive text (\(\hat{b}_{\text{Competence}}\) = 0.50 {[}0.314, 0.676{]}, p \textless{} .001; mean = 3.76, sd = 0.90) than after having read a basic one (mean = 3.26, sd = 0.80). Pooled across all conditions, participants impressiveness ratings were positively associated with competence (\(\hat{b}\) = 0.41 {[}0.312, 0.502{]}, p \textless{} .001).

Participants also trusted a discipline more after having read an impressive text (\(\hat{b}_{\text{trust}}\) = 0.32 {[}0.17, 0.476{]}, p \textless{} .001; mean = 3.68, sd = 0.87) than after having read a basic one. Participants impressiveness ratings were positively associated with trust when pooling across all conditions (\(\hat{b}\) = 0.30 {[}0.206, 0.385{]}, p \textless{} .001).

Participants further had the impression of having learned more after having read the impressive text (mean = 3.66, sd = 0.93; \(\hat{b}\) = 0.89 {[}0.67, 1.108{]}, p \textless{} .001), than after having read the basic one (mean = 2.77, sd = 0.96).

Regarding RQ2, we do, again, not find evidence that consensus modulates the effect of impressiveness hypothesized in H1a and H2a. In Appendix \ref{exp1}, we provide regression tables and figures.

\section{Experiment 2}\label{experiment-2}

In Experiments 1, exposure to impressive scientific content increased trust in science and perceived competence of scientists. In Experiment 2, we wanted to test whether these perceptions were at least partly independent of being able to recall the specific content that induced them. Experiment 2 consisted of a ``main study'' and a ``validation study''.

\subsection{Main study}\label{main-study}

In the main study, after having consented to take part in the study and passing an attention check, participants read the impressive version of one of the vignettes from Experiment 1 (Table \ref{tab:exp1-stimuli}). Whether The text was about archaeology or entomology was randomized. After reading the text, as in Experiment 1, participants were asked about changes in their perception of the scientists' competence (``Would you agree that reading this text has made you think of archaeologists/entomologists as more competent than you thought before?'' {[}1 = Strongly disagree, 2 = Disagree, 3 = Neither agree nor disagree, 4 = Agree, 5 = Strongly agree{]}) and trust in the scientists' discipline (``Having read this text, would you agree that you trust the discipline of archaeology/entomology more than you did before?'' {[}1 - Strongly disagree, 2 - Disagree, 3 - Neither agree nor disagree, 4 - Agree, 5 - Strongly agree{]}). They were also be asked about the impressiveness of the text they read: ``How impressive did you find the findings of the entomologists described in the text you have read?'' (5-point scale: 1 = not impressive at all, 2 = not very impressive, 3 = neither impressive nor unimpressive, 4 = quite impressive, 5 = very impressive). The order of these questions was randomized.

In line with the findings of Experiment 1, we expected participants to have increased perceptions of competence and trust in scientists after:

\textbf{H1a: Participants perceive scientists as more competent after having read an impressive text about their discipline's findings, compared to before.}

\textbf{H1b: Participants trust a discipline more after having read an impressive text about the discipline's findings, compared to before.}

Next, participants were given a recall task: They were asked to rewrite the vignette text they've read just before from memory. They were told that their texts will be read by future participants. To further motivate participants to write down everything they remember, they were also told that they would get a bonus for recalling as much accurate information as possible (without external aids). They were not told how much that bonus would be. We payed them 5p per point gained in the recall task. This way, participants could reach a maximum bonus of 0.8 pound for archaeology (0.05p x 2 points x 8 content elements) and 0.7 pound (0.05p x 2 points x 7 content elements) for entomology.

We predicted that participants would not be able to recall all the information presented in the short vignettes right after having read them (see methods for details):

\textbf{H2: Participants are not able to recall all the information of the original texts.}

This hypothesis, however, only tested whether people forgot any of the content, potentially including non-impressive content. To address this issue, after the recall task, we presented participants with the evaluation grid that we used to asses the open answers from the recall task (see Table \ref{tab:knowledge-evaluation-grid}). For each information element, we asked participants to indicate whether they found it impressive or not (``Do you think this piece of information is impressive?'' {[}Yes; No{]}).

We predicted that even for the subjective subset of information that a participant said they found impressive, they forgot at least some of the content:

\textbf{H3: Participants are not able to recall all the impressive information--as rated by themselves--contained in the original text.}

At the end of the main study, participants were asked about their education level.

\subsection{Validation study}\label{validation-study}

H3 establishes whether participants recall or not content they themselves judge as impressive. However, this judgement could be biased as it happens after the recall task. For a more robust assessment, we conducted a validation study.

The validation study was run XX days after the main study. For this study, a new sample of participants was recruited. After consenting to taking part in the study and passing an attention check, participants were randomly assigned to one of two conditions: In the ``original text'' condition, participants were assigned to read one of the two impressive texts of the main study, again picked at random. In the ``recall text'' condition, participants were also randomly assign to either the archeology or the entomology condition, but instead of the original vignette texts, they read one of the recall texts written by the participants' of the main study. For each participant, a text was randomly sampled (with replacement) from a pool of recall answers. Orthographic and grammatical mistakes in these texts were corrected with the help of ChatGPT beforehand.

We had preregistered to sample among all recall answers from participants of the main study. However, checking recall scores after the main study and before launching the validation study, we found them to be very low on average (REPORT THE AVERAGE, see Appendix - MAKE APPENDIX). To avoid having the answers of ``lazy'' participants--who might have recalled more, but were not motivated to type what they remembered--in our sample, we decided to only selected answers that scored above median in our recall measure for the validation study. This selection is conservative in that it makes it harder to confirm our predictions under H4.

After reading the text, just as in the main study, participants were asked about changes in their perception of the scientists' competence (e.g.~``Would you agree that reading this text has made you think of archaeologists/entomologists as more competent than you thought before?'' {[}1 = Strongly disagree, 2 = Disagree, 3 = Neither agree nor disagree, 4 = Agree, 5 = Strongly agree{]}) and trust in the scientists' discipline (``Having read this text, would you agree that you trust the discipline of archaeology/entomology more than you did before?'' {[}1 - Strongly disagree, 2 - Disagree, 3 - Neither agree nor disagree, 4 - Agree, 5 - Strongly agree{]}). Participants were also be asked about the impressiveness of the text they read: ``How impressive did you find the findings of the entomologists described in the text you have read?'' (5-point scale: 1 = not impressive at all, 2 = not very impressive, 3 = neither impressive nor unimpressive, 4 = quite impressive, 5 = very impressive). The order of these questions was randomized. Finally, participants were asked about their education level.

We predicted that participants in the validation study would be less impressed by the open-ended answers produced by the participants of Experiment 2 compared to the original impressive vignette texts (see Table \ref{tab:exp1-stimuli}) and, accordingly, would have less positively altered perceptions of the scientists' competence and the trustworthiness of their discipline:

\textbf{H4a: The texts produced by participants of the experiment as a result of the recall task will be less impressive than the original texts, as rated by participants of the validation study.}

\textbf{H4b: Participants of the validation study perceive scientists as more competent after having read the original texts, compared to after having read the texts produced by participants of the experiment as a result of the recall task.}

\textbf{H4c: Participants of the validation study trust a discipline more after having read the original texts, compared to after having read the texts produced by participants of the experiment as a result of the recall task.}

\subsection{Methods}\label{methods-1}

\subsubsection{Participants}\label{participants-1}

In our power simulation (see OSF), we varied the sample size of the main study and effect sizes (assuming the same effect size for all hypotheses in each scenario). For all simulations, we assumed a constant validation study sample size of 400 participants (only relevant for H4a, b and c). The power simulation suggested that we would reach a power level of 90\% when assuming a medium effect size of 0.5 with already 50 participants. Due to uncertainty about our assumptions, we recruited a sample 203 participants for the main study, and a sample of 406 participants for the validation study. Although this was not the focus of our simulation, our results showed that for a medium effect size of 0.5, a sample size of 400 for the validation study yielded statistical power of greater than 90\% for all hypotheses based on this sample (H4a, b and c). All participants were from the UK, recruited on the online platform Prolific, and paid to complete the experiment.

We ended up with a final sample of 198 (4 failed attention checks; 99 female, 99 male; \(age_\text{mean}\): 42.65, \(age_\text{sd}\): 15.55, \(age_\text{median}\): 40) for the main study and a sample of 399 (7 failed attention checks; 201 female, 198 male; \(age_\text{mean}\): 41.71, \(age_\text{sd}\): 13.48, \(age_\text{median}\): 40).

\subsubsection{Materials}\label{materials-1}

For Experiment 2, we relied the impressive version of the stimuli used in Experiment 1 (see Table \ref{tab:exp1-stimuli}).

\begin{longtable}[t]{>{\raggedright\arraybackslash}p{20em}>{\raggedright\arraybackslash}p{20em}}
\caption{\label{tab:knowledge-evaluation-grid}Recall evaluation grid}\\
\toprule
Archaeology & Entomology\\
\midrule
1. Archaeologists can determine whether someone was male or female from their bones. & 1. Entomologists use special displays to present images to flies for extremely short periods (less than the blink of an eye).\\
2. Archaeologists can determine how old someone was from their bones. & 2. Entomologists can record how individual cells in flies’ brains react using electrodes.\\
3. Archaeologists can determine whether someone suffered from a range of diseases from their bones. & 3. Entomologists use ultra-precise electron microscopy to examine flies’ eyes.\\
4. Archaeologists can determine at what age someone stopped drinking their mother’s milk, based on the composition of their teeth. & 4. Some flies can perceive images displayed for just three milliseconds. This duration is about ten times shorter than a single movie frame.\\
5. The nerve controlling breathing is larger in humans than in apes. The canal containing that nerve is also larger in humans and Neanderthals than in apes. & 5. Crickets have microscopic hairs situated on antenna-like organs at their rear.\\
\addlinespace
6. The fact that the nerve controlling breathing is larger in humans is possibly due to the need for fine-grained control of breathing to speak. & 6. Crickets' hairs are possibly the most sensitive organs in the animal kingdom. They react to changes in air motion with less energy than one photon.\\
7. Archaeologists determined that most Neanderthals were right-handed, based on analysis of Neanderthals’ tools. & 7. Entomologists measured how cricket hairs react to stimuli, using laser-Doppler velocimetry, which can detect extremely minute movements.\\
8. Handedness is thought to be related to the evolution of language. This suggests that Neanderthals likely possessed a form of language. & \\
\bottomrule
\multicolumn{2}{l}{\rule{0pt}{1em}\textit{Note: }}\\
\multicolumn{2}{l}{\rule{0pt}{1em}For each knowledge element, participants could score a maximum of two points.}\\
\end{longtable}

\paragraph{Recall}\label{recall}

As shown in Table \ref{tab:knowledge-evaluation-grid}, we divided the text into different information elements. For each participant, we calculated a recall score based on how many of these elements they mentioned in their open-ended answer. We coded 0 if a piece of knowledge was not mentioned or was mentioned with very important mistakes (e.g., writing ``the nerve controlling fine hand movement is bigger in humans'' instead of ``the nerve controlling breathing is bigger in humans''). We coded 1 if the piece of knowledge was mentioned, but some important elements were missing (e.g., writing ``Archaeologists can determine at what age someone stopped drinking their mother's milk'' instead of « Archaeologists can determine at what age someone stopped drinking their mother's milk from the composition of their teeth »), and/or there were some mistakes (e.g., writing ``Archaeologists can determine at what age someone stopped drinking their mother's milk, based on the bones'' instead of ``Archaeologists can determine at what age someone stopped drinking their mother's milk based on the teeth''). We coded 2 if the piece of knowledge was mentioned with all the main content, even if the participant had not used the precise technical words (e.g., ``neanderthals'', ``laser-Doppler velocimetry'') or had changed the phrasing of the information in other ways. We coded participants' answers with the help of ChatGPT. We provided ChatGPT with the evaluation grid and the following prompt:

\begin{quote}
\textbf{Prompt ChatGPT}: \emph{``In an experiment, participants read the following vignette : {[}\ldots{]}. Then, participants are asked to write what they recall from the vignette. Evaluate whether the following piece of information is included in the participant's text : {[}\ldots{]}. You should give one of three grades : 0 = the piece of knowledge is not mentioned, or it is mentioned with very important mistakes (e.g.~writing ``the nerve controlling fine hand movement is bigger in humans'' instead of ``the nerve controlling breathing is bigger in humans''). 1 = the piece of knowledge is mentioned, but some important elements are missing (e.g.~writing ``Archaeologists can determine at what age someone stopped drinking their mother's milk'' instead of « Archaeologists can determine at what age someone stopped drinking their mother's milk from the composition of their teeth »), and/or there are some mistakes (e.g.~writing ``Archaeologists can determine at what age someone stopped drinking their mother's milk, based on the bones'' instead of ``Archaeologists can determine at what age someone stopped drinking their mother's milk based on the teeth''). 2 = the piece of knowledge is mentioned with all the main elements. Note: a grade of 2 is attributed even if the participant has not used the precise technical words (e.g.~``neanderthals'', ``laser-Doppler velocimetry'') or has changed the phrasing of the information in other ways. Provide a brief explanation, and conclude the explanation with the score formatted as Score= followed by the numerical value, with no text or symbol after the score. Here is the participant's text: {[}\ldots{]}.''}
\end{quote}

To validate the scores assigned by ChatGPT, we compared them to scores assigned by two human coders for a subsample of 80 randomly chosen texts (half on archaeology, half on entomology). The human coders were unaware of the study context and the hypotheses. They read the same prompt provided to ChatGPT as their instructions. They received minimal training on some texts obtained in a pilot study. To measure the reliability between ChatGPT coding and human coding, we calculated an ICC of inter-rater agreement for two-way designs with random raters (\textbf{tenhoveUpdatedGuidelinesSelecting2024?}). Following the guidelines of (\textbf{heymanBehavioralObservationCoding2014?}), we took 0.7 as a threshold for acceptable reliability. If we met this threshold, we used only ChatGPT evaluations for constructing the recall scores. If we did not, we qualitatively investigated differences in decisions between human coders and ChatGPT. After this assessment, we might have tried using a different prompt for ChatGPT and/or adding other human coders.

Since the two vignettes contained a slightly different number of total information elements according to our evaluation grid (8 for archaeology, 7 for entomology), we used a relative measure for the final recall score, namely the share of obtained points among all possible points.

\paragraph{Recall of impressive items}\label{recall-of-impressive-items}

After giving their post-evaluation of scientists' competence and trust in the different disciplines, participants were presented with the evaluation grid shown in Table \ref{tab:knowledge-evaluation-grid} for the respective discipline they had been randomized to see. For each element in the evaluation grid, they were asked whether they found it impressive or not. For each participant, we computed the recall score they obtained on those elements they subjectively rated as impressive. Since the number of items seen as impressive varied within participants, the final recall measure for impressive elements was a relative one, namely the share of points among all possible points for the items rated as impressive by the participant.

\subsection{Results and discussion}\label{results-and-discussion-1}

\subsubsection{Main study}\label{main-study-1}

For the first set of hypotheses, we tested whether the average scores for perceived change in competence and trust are significantly different from their respective scale midpoint (which corresponds to no change in perception). Practically, we subtracted the scale mid-point from from all answers to then test if the resulting scores are statistically different from 0. We first tested the outcome variables' distributions for normality, using a Shapiro-Wilk test. In all cases, this test suggested that the data is considered non-normally distributed (p \textless{} 0.05). Following our pre-registration, we therefor did not run a default one-sample t-test, but used a Wilcoxon signed-rank test instead. This test is a non-parametric alternative to the one-sample t-test, used when the normality assumption is violated. It tests if the median of the outcome distributions is significantly different from 0.

We confirm both H1a and H1b, finding that having read an impressive text about the findings of a scientific discipline, participants saw the scientists as more competent (median = 4, W = 12,388.50, p \textless{} .001). To test H2 and H3, we deviated from our pre-registration and select only participants with an above median recall score. This selection is because we realized many very low recall scores. To avoid having ``lazy'' participants bias our measure in direction of our predictions, we decided for this conservative approach.

We tested whether the recall score for the best 50\% of participants was statistically different from perfect recall. Practically, we tested whether the average forgetting score (1-recall score, since recall score is a percentage) differs significantly from zero. Since a Shapiro-Wilk test suggested a non-normal population distribution, we followed our preregistration in using a Wilcoxon signed-rank test. Confirming both hypotheses, the test provides evidence that participants did not perfectly recall all information (median = 0.43, W = 5778, p \textless{} .001) and they did not recall all information contained in the elements they judged as impressive themselves after the taks (median = 0.43, W = 5778, p \textless{} .001)\footnote{Both median values are not the same, but highly similar, because participants rated most knowledge elements as impressive, making the general recall score and the recall score for impressive elements very similar}. Note that most participants declared finding all knowledge elements impressive (\(mean_{\text{Archeology}}\) = 7.04, \(median_{\text{Archeology}}\) = 8, number of knowledge elements = 8; \(mean_{\text{Entomology}}\) = 6.39, \(median_{\text{Entomology}}\) = 7, number of knowledge elements = 7).

\subsubsection{Validation study}\label{validation-study-1}

For H4a, b and c, we compared trust, competence and impressiveness ratings between the ``original texts'' and the ``recall texts'' conditions of the validation study using independent sample t-tests. Participants who read the original texts reported being more impressed (\(\hat{b}_{\text{Impressiveness}}\) = 0.39, t = 5.41, p \textless{} .001), rated the scientists of the respective discipline as more competent (\(\hat{b}_{\text{Confidence}}\) = 0.27, t = 3.07, p p = 0.002) and had more trust in the respective discipline (\(\hat{b}_{\text{Trust}}\) = 0.27, t = 3.40, p \textless{} .001) than participants who read the texts of the recall task.

\section{Overall discussion}\label{overall-discussion}

The rational impression account of trust in science predicts that people trust science primarily because they are impressed by it, not because they remember much explicit knowledge. This account is rational, in that the heuristics it builds on are sound in many contexts. Trusting without understanding is not irrational: If someone finds out something that is hard to know (and consensual, but here we take that for granted), such as the distance of the milky way, we should expect them to be competent, even without knowing the details. Forgetting specific content is not irrational either: It has been argued that episodic memory exits for us to be able to justify our beliefs in communication with others (Mahr \& Csibra, 2018). As a consequence, we should be particularly good at remembering things we might need to convince others. There is no reason for young learners' minds to anticipate having to convince others of the merits of science: At school, where most people have most of their encounters with science, everyone learns the same thing and accepts it.

In two experiments, we have provided evidence for this account. In Experiment 1, we have shown that impressive scientific findings lead people to think of scientists as more competent and trust science more. In Experiment 2, we have shown that these impressions persist despite participants forgetting specific content that generated them almost immediately after reading it.

It has long been a puzzle to the deficit model--the paradigmatic explanation for trust in science which assumes that it is primarily the result of knowledge about science--that across many studies, knowledge about science is at best weakly associated with science attitudes, while the association between trust in science and education, and in particular science education, is strong. The rational impression account can make sense of this. It suggests that the mechanism by which education shapes trust is the exposure to impressive science, rather than the creation of long-lasting specific knowledge or even profound understanding.

Are we suggesting that the goal of education to foster understanding of science is misguided? No.~For science to be impressive, future scientists must have understood the results of their predecessors. And, arguably, a better understanding how difficult it was to obtain a certain result should make science even more impressive. Our account generally stresses the importance of science communication: Well-crafted science communication highlighting the feat of scientists can make a difference in shaping trust in science. Because people struggle to remember what exactly impressed them, we cannot expect trust in science to automatically diffuse in the population from hearsay. Someone trusting medicine might find it hard to convince someone skeptical of vaccines, because trust for them is an abstract impression largely disconnected from specific medical knowledge. Therefore, science communication needs to be actively pursued by scientists, policy makers, and educators.

The present experiments have a number of limitations: First, they were run on convenience samples recruited in a single country, the UK. Second, they were run within a very short time frame. While we can show that participants almost immediately forget about impressive content, it is not clear from our study for how long the impressions persists. We hope that future studies will extend our findings to other populations and longer time frames.

\FloatBarrier

\section{References}\label{references}

\phantomsection\label{refs}
\begin{CSLReferences}{1}{0}
\bibitem[\citeproctext]{ref-alganTrustScientistsTimes2021}
Algan, Y., Cohen, D., Davoine, E., Foucault, M., \& Stantcheva, S. (2021). Trust in scientists in times of pandemic: Panel evidence from 12 countries. \emph{Proceedings of the National Academy of Sciences}, \emph{118}(40), e2108576118. \url{https://doi.org/10.1073/pnas.2108576118}

\bibitem[\citeproctext]{ref-allumScienceKnowledgeAttitudes2008}
Allum, N., Sturgis, P., Tabourazi, D., \& Brunton-Smith, I. (2008). Science knowledge and attitudes across cultures: a meta-analysis. \emph{Public Understanding of Science}, \emph{17}(1), 35--54. \url{https://doi.org/10.1177/0963662506070159}

\bibitem[\citeproctext]{ref-bakEducationPublicAttitudes2001}
Bak, H.-J. (2001). Education and Public Attitudes toward Science: Implications for the {``}Deficit Model{''} of Education and Support for Science and Technology. \emph{Social Science Quarterly}, \emph{82}(4), 779--795. \url{https://doi.org/10.1111/0038-4941.00059}

\bibitem[\citeproctext]{ref-bauerWhatCanWe2007}
Bauer, M. W., Allum, N., \& Miller, S. (2007). What can we learn from 25 years of PUS survey research? Liberating and expanding the agenda. \emph{Public Understanding of Science}, \emph{16}(1), 79--95. \url{https://doi.org/10.1177/0963662506071287}

\bibitem[\citeproctext]{ref-beckRiskSocietyNew1992}
Beck, U. (1992). \emph{Risk society: towards a new modernity} (Repr). London: Sage.

\bibitem[\citeproctext]{ref-colognaRoleTrustClimate2020}
Cologna, V., \& Siegrist, M. (2020). The role of trust for climate change mitigation and adaptation behaviour: A meta-analysis. \emph{Journal of Environmental Psychology}, \emph{69}, 101428. \url{https://doi.org/10.1016/j.jenvp.2020.101428}

\bibitem[\citeproctext]{ref-durantPublicUnderstandingScience1989}
Durant, J. R., Evans, G. A., \& Thomas, G. P. (1989). The public understanding of science. \emph{Nature}, \emph{340}(6228), 11--14. \url{https://doi.org/10.1038/340011a0}

\bibitem[\citeproctext]{ref-feinsteinSustainedExperienceEmotion2010}
Feinstein, J. S., Duff, M. C., \& Tranel, D. (2010). Sustained experience of emotion after loss of memory in patients with amnesia. \emph{Proceedings of the National Academy of Sciences}, \emph{107}(17), 7674--7679. \url{https://doi.org/10.1073/pnas.0914054107}

\bibitem[\citeproctext]{ref-funkPublicConfidenceScientists2020}
Funk, C., \& Kennedy, B. (2020). \emph{Public confidence in scientists has remained stable for decades}.

\bibitem[\citeproctext]{ref-funkScienceScientistsHeld2020}
Funk, C., Tyson, A., Kennedy, B., \& Johnson, C. (2020). \emph{Science and scientists held in high esteem across global publics}. Retrieved from \url{https://www.pewresearch.org/science/2020/09/29/science-and-scientists-held-in-high-esteem-across-global-publics/}

\bibitem[\citeproctext]{ref-gauchatCulturalAuthorityScience2011}
Gauchat, G. (2011). The cultural authority of science: Public trust and acceptance of organized science. \emph{Public Understanding of Science}, \emph{20}(6), 751--770. \url{https://doi.org/10.1177/0963662510365246}

\bibitem[\citeproctext]{ref-gauchatPoliticizationSciencePublic2012}
Gauchat, G. (2012). Politicization of Science in the Public Sphere: A Study of Public Trust in the United States, 1974 to 2010. \emph{American Sociological Review}, \emph{77}(2), 167--187. \url{https://doi.org/10.1177/0003122412438225}

\bibitem[\citeproctext]{ref-gauchatCulturalCognitiveMappingScientific2018}
Gauchat, G., \& Andrews, K. T. (2018). The Cultural-Cognitive Mapping of Scientific Professions. \emph{American Sociological Review}, \emph{83}(3), 567--595. \url{https://doi.org/10.1177/0003122418773353}

\bibitem[\citeproctext]{ref-giddensModernitySelfidentitySelf1991}
Giddens, A. (1991). \emph{Modernity and self-identity: Self and society in the late modern age}. Stanford, Calif: Stanford University Press.

\bibitem[\citeproctext]{ref-habermasJurgenHabermasSociety1989}
Habermas, J. (1989). \emph{Jürgen Habermas on society and politics: a reader} (1. {[}print.{]}; S. Seidman, Ed.). Boston: Beacon Press.

\bibitem[\citeproctext]{ref-lindholtPublicAcceptanceCOVID192021}
Lindholt, M. F., Jørgensen, F., Bor, A., \& Petersen, M. B. (2021). Public acceptance of COVID-19 vaccines: cross-national evidence on levels and individual-level predictors using observational data. \emph{BMJ Open}, \emph{11}(6), e048172. \url{https://doi.org/10.1136/bmjopen-2020-048172}

\bibitem[\citeproctext]{ref-liquinMotivatedLearnAccount2022}
Liquin, E. G., \& Lombrozo, T. (2022). Motivated to learn: An account of explanatory satisfaction. \emph{Cognitive Psychology}, \emph{132}, 101453. \url{https://doi.org/10.1016/j.cogpsych.2021.101453}

\bibitem[\citeproctext]{ref-mahrWhyWeRemember2018}
Mahr, J. B., \& Csibra, G. (2018). Why do we remember? The communicative function of episodic memory. \emph{Behavioral and Brain Sciences}, \emph{41}, e1. \url{https://doi.org/10.1017/S0140525X17000012}

\bibitem[\citeproctext]{ref-millerMeasurementCivicScientific1998a}
Miller, J. D. (1998). The measurement of civic scientific literacy. \emph{Public Understanding of Science}, \emph{7}(3), 203--223. \url{https://doi.org/10.1088/0963-6625/7/3/001}

\bibitem[\citeproctext]{ref-millerPublicUnderstandingAttitudes2004}
Miller, J. D. (2004). Public Understanding of, and Attitudes toward, Scientific Research: What We Know and What We Need to Know. \emph{Public Understanding of Science}, \emph{13}(3), 273--294. \url{https://doi.org/10.1177/0963662504044908}

\bibitem[\citeproctext]{ref-nationalacademiesofsciencesengineeringandmedicineScienceLiteracyConcepts2016}
National Academies of Sciences, Engineering, and Medicine. (2016). \emph{Science Literacy: Concepts, Contexts, and Consequences} (C. E. Snow \& K. A. Dibner, Eds.). Washington, D.C.: National Academies Press. \url{https://doi.org/10.17226/23595}

\bibitem[\citeproctext]{ref-noyScienceGoodEffects2019}
Noy, S., \& O'Brien, T. L. (2019). Science for good? The effects of education and national context on perceptions of science. \emph{Public Understanding of Science}, \emph{28}(8), 897--916. \url{https://doi.org/10.1177/0963662519863575}

\bibitem[\citeproctext]{ref-pfanderHowWiseCrowd2025}
Pfänder, J., De Courson, B., \& Mercier, H. (2025). How wise is the crowd: Can we infer people are accurate and competent merely because they agree with each other? \emph{Cognition}, \emph{255}, 106005. \url{https://doi.org/10.1016/j.cognition.2024.106005}

\bibitem[\citeproctext]{ref-pfanderQuasiuniversalAcceptanceBasic2024}
Pfänder, J., Kerzreho, L., \& Mercier, H. (2024). \emph{Quasi-universal acceptance of basic science in the US}. \url{https://doi.org/10.31219/osf.io/qc43v}

\bibitem[\citeproctext]{ref-smithTrendsPublicAttitudes2013}
Smith, T. W., \& Son, J. (2013). Trends in public attitudes about confidence in institutions. \emph{General Social Survey Final Report}.

\bibitem[\citeproctext]{ref-sturgisTrustScienceSocial2021}
Sturgis, P., Brunton-Smith, I., \& Jackson, J. (2021). Trust in science, social consensus and vaccine confidence. \emph{Nature Human Behaviour}, \emph{5}(11), 1528--1534. \url{https://doi.org/10.1038/s41562-021-01115-7}

\bibitem[\citeproctext]{ref-wellcomeglobalmonitorWellcomeGlobalMonitor2018}
Wellcome Global Monitor. (2018). \emph{Wellcome Global Monitor 2018}. Retrieved from \url{https://wellcome.org/reports/wellcome-global-monitor/2018}

\bibitem[\citeproctext]{ref-wellcomeglobalmonitorWellcomeGlobalMonitor2020}
Wellcome Global Monitor. (2020). \emph{Wellcome Global Monitor 2020: Covid-19}. Retrieved from \url{https://wellcome.org/reports/wellcome-global-monitor-covid-19/2020}

\end{CSLReferences}

\newpage

\appendix



\end{document}
